% appendix.tex — generated by building_website/scripts/build_appendix.py.
% Re-running the script preserves whatever sits between each row's
% `% NOTES-BEGIN <ref>` / `% NOTES-END <ref>` markers, so hand-written
% notes are never overwritten when panel data is refreshed.
%
% Preamble (documentclass, packages, custom macros, minted setup) is
% supplied by the wrapping document; this file only emits body content.

%% ============================================================

\section*{Section 3.1}

%% ROW-BEGIN def_3_1
\subsection*{Definition: Conditional directed mixed graphs (CDMG)}

\noindent
\begin{minipage}[t]{0.49\textwidth}
\begin{Def}[Conditional directed mixed graphs (CDMG)]
    \label{def-cdmg}
    A \emph{conditional directed mixed graph (CDMG)} $G$---per definition---consists of two finite, disjoint sets of
    vertices (also called nodes):
    \begin{enumerate}[label=\roman*.)]
        \item $J$, whose elements are called input nodes,
        \item $V$, whose elements are called output nodes,
    \end{enumerate}
    with both $J$ and $V$ finite and $J \cap V = \emptyset$;
%    such that $J \cup V \neq \emptyset$;
    and two sets of edges (here ``disjoint'' is to be read as trivial type-level disjointness only: directed edges and bidirected edges are objects of distinct types, and no graph-theoretic mutual-exclusion is imposed between $E$ and $L$):
    \begin{enumerate}[resume,label=\roman*.)]
        \item $E \subseteq (J \cup V) \times V$, the set of directed edges,
        \item $L \subseteq V \times V$, the set of bidirected edges, subject to the irreflexivity and symmetry constraint
            \[ \forall\, (v_1, v_2) \in V \times V :\quad (v_1, v_2) \in L \;\Longrightarrow\; v_1 \neq v_2 \,\land\, (v_2, v_1) \in L. \]
            Equivalently, $L$ may be identified with a subset of the quotient $(V \times V)/{\sim}$ under the swap relation $(v_1, v_2) \sim (v_2, v_1)$ (i.e.\ a set of unordered pairs of distinct elements of $V$); under the ordered-pair encoding above, membership of $(v_1, v_2)$ and of $(v_2, v_1)$ in $L$ always coincide, and the irreflexivity condition $v_1 \neq v_2$ applies under either encoding.
        \item[] In particular, for $v_1, v_2 \in V$ with $v_1 \neq v_2$, it is permitted to have $(v_1, v_2) \in E$ and the unordered pair $\{v_1, v_2\}$ in $L$ simultaneously; under the ordered-pair encoding of $L$, the same ordered pair $(v, w) \in V \times V$ may belong to both $E$ and $L$.
    \end{enumerate}
\end{Def}
\end{minipage}\hfill
\begin{minipage}[t]{0.49\textwidth}
\begin{minted}{lean4}
structure CDMG (Node : Type*) [DecidableEq Node] where
  J : Finset Node
  V : Finset Node
  hJV_disj : Disjoint J V
  E : Finset (Node × Node)
  hE_subset : ∀ ⦃e : Node × Node⦄, e ∈ E → e.1 ∈ J ∪ V ∧ e.2 ∈ V
  L : Finset (Node × Node)
  hL_subset : ∀ ⦃e : Node × Node⦄, e ∈ L → e.1 ∈ V ∧ e.2 ∈ V
  hL_irrefl : ∀ ⦃v1 v2 : Node⦄, (v1, v2) ∈ L → v1 ≠ v2
  hL_symm : ∀ ⦃v1 v2 : Node⦄, (v1, v2) ∈ L → (v2, v1) ∈ L
\end{minted}
\end{minipage}

\vspace{0.5em}
\noindent\textbf{Notes:}\par
% NOTES-BEGIN def_3_1

% NOTES-END def_3_1
\vspace{5em}
\noindent\hrulefill
\par\bigskip
%% ROW-END def_3_1

%% ROW-BEGIN def_3_2
\subsection*{Definition: CDMGNotation}

\noindent
\begin{minipage}[t]{0.49\textwidth}
\begin{Not}[CDMGNotation]
    \label{not-cdmg}
    Let $G = (J, V, E, L)$ be a CDMG (def \ref{def-cdmg}).
    We introduce the following seven shorthand notations, each defined as a literal abbreviation for a set-theoretic statement over the four components $J, V, E, L$ of $G$.
    The visual macros $\tuh, \hut, \huh, \suh, \hus, \sus$ are notation only; the load-bearing definitional content is the right-hand side of each abbreviation below.
    \begin{enumerate}
        %\item $G(V|\doit(J))$ to represent the graph $(J,V,E,L)$, where $E$ and $L$ are kept implicit in this notation.
        \item $v \in G \;:\Longleftrightarrow\; v \in J \cup V$.
        \item $v_1 \tuh v_2 \in G \;:\Longleftrightarrow\; (v_1, v_2) \in E$.
        \item $v_1 \hut v_2 \in G \;:\Longleftrightarrow\; (v_2, v_1) \in E$.
        \item $v_1 \huh v_2 \in G \;:\Longleftrightarrow\; (v_1, v_2) \in L$.
              Because $L$ is required to be symmetric by the CDMG definition (def \ref{def-cdmg}: $(v_1, v_2) \in L \Longleftrightarrow (v_2, v_1) \in L$ for all $v_1, v_2$), this abbreviation is itself symmetric in its two arguments:
              \[ v_1 \huh v_2 \in G \;\Longleftrightarrow\; v_2 \huh v_1 \in G. \]
        \item $v_1 \suh v_2 \in G \;:\Longleftrightarrow\; \bigl(v_1 \tuh v_2 \in G\bigr) \,\lor\, \bigl(v_1 \huh v_2 \in G\bigr)$, equivalently $(v_1, v_2) \in E \,\lor\, (v_1, v_2) \in L$.
              Because the $\huh$-disjunct is symmetric in its arguments (item~4), swapping $v_1 \leftrightarrow v_2$ in this disjunction turns $\suh$ into $\hus$:
              \[ v_1 \suh v_2 \in G \;\Longleftrightarrow\; v_2 \hus v_1 \in G. \]
        \item $v_1 \hus v_2 \in G \;:\Longleftrightarrow\; \bigl(v_1 \hut v_2 \in G\bigr) \,\lor\, \bigl(v_1 \huh v_2 \in G\bigr)$, equivalently $(v_2, v_1) \in E \,\lor\, (v_1, v_2) \in L$.
              By the same $\huh$-symmetry of item~4, $\hus$ becomes $\suh$ under swap of its arguments:
              \[ v_1 \hus v_2 \in G \;\Longleftrightarrow\; v_2 \suh v_1 \in G. \]
        \item $v_1 \sus v_2 \in G \;:\Longleftrightarrow\; \bigl(v_1 \tuh v_2 \in G\bigr) \,\lor\, \bigl(v_1 \hut v_2 \in G\bigr) \,\lor\, \bigl(v_1 \huh v_2 \in G\bigr)$, equivalently $(v_1, v_2) \in E \,\lor\, (v_2, v_1) \in E \,\lor\, (v_1, v_2) \in L$.

              \emph{Exhaustiveness of the three disjuncts.}
              These three configurations enumerate every edge configuration a CDMG admits between an ordered pair of nodes $(v_1, v_2)$.
              Although mechanically reading the placeholder convention ``$\star$ stands for arrowhead or tail'' would suggest a fourth, tail-tail case ``$v_1 \tut v_2$'', no edge type of a CDMG carries tails on both ends: per def \ref{def-cdmg}, directed edges in $E$ have one tail and one head, and bidirected edges in $L$ have two heads, so no tail-tail edge type exists.
              The case $v_1 \tut v_2$ is therefore vacuous as a CDMG edge and is excluded by definition, not omitted by oversight.
              Any case analysis on ``$v_1 \sus v_2 \in G$'' may assume exhaustiveness over exactly the three listed alternatives $\{\tuh, \hut, \huh\}$.

              \emph{Symmetry of $\sus$.}
              Under the swap $v_1 \leftrightarrow v_2$, the three disjuncts pair off: $v_1 \tuh v_2 \in G$ exchanges with $v_2 \hut v_1 \in G$ (items~2~$\leftrightarrow$~3), and $v_1 \huh v_2 \in G$ exchanges with $v_2 \huh v_1 \in G$ by the $\huh$-symmetry of item~4. Hence $\sus$ is symmetric in its two arguments:
              \[ v_1 \sus v_2 \in G \;\Longleftrightarrow\; v_2 \sus v_1 \in G. \]
    \end{enumerate}
    A ``$\star$'' appearing in the diagrammatic shorthands ($\suh$, $\hus$, $\sus$) is a placeholder reading ``arrowhead or tail''; in $\sus$ the would-be tail-tail combination is vacuous as a CDMG edge type (see item~7), so $\sus$ ranges over exactly the three configurations enumerated there.
\end{Not}
\end{minipage}\hfill
\begin{minipage}[t]{0.49\textwidth}
\begin{minted}{lean4}
variable {Node : Type*} [DecidableEq Node]

instance instMembership : Membership Node (CDMG Node) where
  mem G v := v ∈ G.J ∪ G.V

def tuh (G : CDMG Node) (v1 v2 : Node) : Prop := (v1, v2) ∈ G.E

def hut (G : CDMG Node) (v1 v2 : Node) : Prop := (v2, v1) ∈ G.E

def huh (G : CDMG Node) (v1 v2 : Node) : Prop := (v1, v2) ∈ G.L

def suh (G : CDMG Node) (v1 v2 : Node) : Prop :=
  G.tuh v1 v2 ∨ G.huh v1 v2

def hus (G : CDMG Node) (v1 v2 : Node) : Prop :=
  G.hut v1 v2 ∨ G.huh v1 v2

def sus (G : CDMG Node) (v1 v2 : Node) : Prop :=
  G.tuh v1 v2 ∨ G.hut v1 v2 ∨ G.huh v1 v2
\end{minted}
\end{minipage}

\vspace{0.5em}
\noindent\textbf{Notes:}\par
% NOTES-BEGIN def_3_2

% NOTES-END def_3_2
\vspace{5em}
\noindent\hrulefill
\par\bigskip
%% ROW-END def_3_2

%% ROW-BEGIN def_3_3
\subsection*{Definition: EdgeRelations}

\noindent
\begin{minipage}[t]{0.49\textwidth}
\begin{Def}[EdgeRelations]
    \label{def-edge-relations}
    Let $G = (J, V, E, L)$ be a CDMG (def \ref{def-cdmg}); throughout we use the notation of def \ref{not-cdmg}, and recall that $E \subseteq (J \cup V) \times V$ and $L \subseteq V \times V$, with $L$ symmetric in its two coordinates.

    \begin{enumerate}[label=\roman*.)]
        \item \emph{Adjacency.}
              For $v_1, v_2 \in J \cup V$, the nodes $v_1$ and $v_2$ are called \emph{adjacent in $G$} iff
              \[ (v_1, v_2) \in E \;\lor\; (v_2, v_1) \in E \;\lor\; (v_1, v_2) \in L, \]
              equivalently iff $v_1 \sus v_2 \in G$ (def \ref{not-cdmg}, item~7).

        \item \emph{Edge into a vertex.}
              The relations \emph{into} and \emph{out of} below classify \emph{edges} of $G$ (the underlying ordered pairs of $E$ or of $L$) relative to a single bound vertex parameter $v$ -- not relative to a particular labelling $v_1, v_2$ of the edge's endpoints.
              For an ordered pair $e = (e_1, e_2) \in (J \cup V) \times (J \cup V)$ and a vertex $v \in J \cup V$, we call $e$ an \emph{edge into $v$ (in $G$)} iff
              \[
                \bigl(e \in E \;\land\; e_2 = v\bigr) \;\lor\; \bigl(e \in L \;\land\; (e_1 = v \,\lor\, e_2 = v)\bigr).
              \]
              In particular, only ordered pairs $e \in E \cup L$ can be edges into $v$.

        \item \emph{Edge out of a vertex.}
              For an ordered pair $e = (e_1, e_2) \in (J \cup V) \times (J \cup V)$ and a vertex $v \in J \cup V$, we call $e$ an \emph{edge out of $v$ (in $G$)} iff
              \[ e \in E \;\land\; e_1 = v. \]
              In particular, no $L$-edge is out of any vertex.

        \item \emph{Reconciliation with the source-block writings $v_1 \tuh v_2$, $v_1 \hut v_2$, $v_1 \huh v_2$.}
              The two writings $v_1 \tuh v_2$ and $v_2 \hut v_1$ both abbreviate the single underlying statement $(v_1, v_2) \in E$ (def \ref{not-cdmg}, items~2--3); and the two writings $v_1 \huh v_2$ and $v_2 \huh v_1$ both abbreviate $(v_1, v_2) \in L$, which moreover coincides with $(v_2, v_1) \in L$ by the symmetry of $L$ (def \ref{def-cdmg}).
              Hence membership of an edge in the categories \emph{into $v$} or \emph{out of $v$} depends only on the underlying element of $E$ or of $L$, not on which of these textual writings is used; in particular, the source block's listing $v_1 \tuh v_2$ \emph{or} $v_2 \hut v_1$ in its item~3 and its single-writing listings in item~2 are equivalent under this identification.

              Spelled out per writing, items~ii.--iii.\ above agree with the source block's items~2--3 as follows.
              \begin{itemize}
                  \item An $E$-edge written $v_1 \tuh v_2$ -- i.e.\ underlying ordered pair $(v_1, v_2) \in E$, equivalently written $v_2 \hut v_1$ -- is \emph{into $v_2$} (clause ii.\ fires via $e \in E$ with $e_2 = v_2$) and \emph{out of $v_1$} (clause iii.\ fires via $e \in E$ with $e_1 = v_1$).
                  \item An $E$-edge written $v_1 \hut v_2$ -- i.e.\ underlying ordered pair $(v_2, v_1) \in E$, equivalently written $v_2 \tuh v_1$ -- is \emph{into $v_1$} and \emph{out of $v_2$}, by the same clauses with $v_1$ and $v_2$ interchanged.
                  \item An $L$-edge written $v_1 \huh v_2$ -- i.e.\ underlying ordered pair $(v_1, v_2) \in L$, equivalently $(v_2, v_1) \in L$ and equivalently written $v_2 \huh v_1$ -- is \emph{into $v_1$ and into $v_2$} (the $L$-clause of ii.\ fires for either endpoint), and is \emph{out of neither} $v_1$ nor $v_2$ (clause iii.\ requires $e \in E$, whereas $e \in L$).
              \end{itemize}
              In particular, \emph{into $v$} and \emph{out of $v$} are not complementary classifications of edges incident at $v$: an $L$-edge incident at $v$ is into $v$ but never out of $v$.
    \end{enumerate}
\end{Def}
\end{minipage}\hfill
\begin{minipage}[t]{0.49\textwidth}
\begin{minted}{lean4}
variable {Node : Type*} [DecidableEq Node]

def adjacent (G : CDMG Node) (v1 v2 : Node) : Prop := G.sus v1 v2

def into (G : CDMG Node) (v : Node) (e : Node × Node) : Prop :=
  (e ∈ G.E ∧ e.2 = v) ∨ (e ∈ G.L ∧ (e.1 = v ∨ e.2 = v))

def outOf (G : CDMG Node) (v : Node) (e : Node × Node) : Prop :=
  e ∈ G.E ∧ e.1 = v
\end{minted}
\end{minipage}

\vspace{0.5em}
\noindent\textbf{Notes:}\par
% NOTES-BEGIN def_3_3

% NOTES-END def_3_3
\vspace{5em}
\noindent\hrulefill
\par\bigskip
%% ROW-END def_3_3

%% ROW-BEGIN claim_3_1
\subsection*{Claim: CDMGRestrictions}

\noindent
\begin{minipage}[t]{0.49\textwidth}
% --- statement (restated) ---------------------------------------------------
\begin{Rem}[CDMGRestrictions]
    Let $G = (J, V, E, L)$ be a CDMG (def \ref{def-cdmg}); we use the notation of \ref{not-cdmg} throughout. Recall, per def \ref{def-cdmg}, that $E \subseteq (J \cup V) \times V$, that $L \subseteq V \times V$ with $L$ symmetric and irreflexive, and that $J \cap V = \emptyset$.

    The restrictions encoded in def \ref{def-cdmg} have the following three consequences. In the LN wording, the clauses ``$j \hus v \notin G$'' and ``edges $j \tuh v$ are allowed'' carry an implicit universal quantification over the second-position variable $v$; we surface these quantifiers explicitly below, with $v$ ranging over $J \cup V$ in the first clause and over $V$ only in the second (see also the addition spelled out under sub-claim ii.).

    \begin{enumerate}[label=\roman*.)]
        \item \emph{No edge with an arrowhead at a $J$-node.}
              For every $j \in J$ and every $v \in J \cup V$,
              \[ j \hus v \notin G. \]
              Unfolding the shorthand $\hus$ via def \ref{not-cdmg} item~6 (so $j \hus v \in G$ abbreviates $(v, j) \in E \lor (j, v) \in L$), this is equivalent to the set-theoretic conjunction
              \[ (v, j) \notin E \;\land\; (j, v) \notin L. \]

        \item \emph{The type of $G$ does not forbid a directed edge from a $J$-node into a $V$-node.}
              For every $j \in J$ and every $v \in V$, the ordered pair $(j, v)$ satisfies $j \in J \cup V$ and $v \in V$, so the typing constraint $E \subseteq (J \cup V) \times V$ of def \ref{def-cdmg} is compatible with $(j, v) \in E$; equivalently, no restriction in def \ref{def-cdmg} forbids the inclusion of $(j, v)$ in $E$ (in the notation of \ref{not-cdmg}: no restriction forbids $j \tuh v \in G$). This is a \emph{permission} statement, not an existence statement: $(j, v) \in E$ is admissible by the definition for every $(j, v) \in J \times V$, but is not asserted to hold.

        \item \emph{No two $J$-nodes are adjacent.}
              For every $j_1 \in J$ and every $j_2 \in J$, the nodes $j_1$ and $j_2$ are \emph{not} adjacent in $G$ in the sense of def \ref{def-edge-relations} item~i.; equivalently,
              \[ (j_1, j_2) \notin E \;\land\; (j_2, j_1) \notin E \;\land\; (j_1, j_2) \notin L. \]
    \end{enumerate}
\end{Rem}

% --- proof ------------------------------------------------------------------
\begin{proof}
All three sub-claims rest on the same underlying tool: the typing constraints carried by def \ref{def-cdmg}, namely $E \subseteq (J \cup V) \times V$, $L \subseteq V \times V$, and $J \cap V = \emptyset$. We prove each in turn.

\medskip
\emph{Proof of (i).}
Fix $j \in J$ and $v \in J \cup V$, and suppose, for contradiction, that $j \hus v \in G$. By def \ref{not-cdmg} item~6, the shorthand $j \hus v \in G$ abbreviates the disjunction
\[ (v, j) \in E \;\lor\; (j, v) \in L. \]
We rule out both disjuncts.
\begin{itemize}
    \item If $(v, j) \in E$, then by the typing constraint $E \subseteq (J \cup V) \times V$ of def \ref{def-cdmg}, the second component of any element of $E$ lies in $V$; hence $j \in V$. But $j \in J$ by assumption and $J \cap V = \emptyset$ (def \ref{def-cdmg}), so $j \in J \cap V = \emptyset$, a contradiction.
    \item If $(j, v) \in L$, then by the typing constraint $L \subseteq V \times V$ of def \ref{def-cdmg}, the first component of any element of $L$ lies in $V$; hence $j \in V$. The same contradiction $j \in J \cap V = \emptyset$ follows.
\end{itemize}
Both disjuncts therefore fail, so $j \hus v \notin G$. Unfolding the shorthand once more, this is the conjunction $(v, j) \notin E \land (j, v) \notin L$, as stated.

\medskip
\emph{Proof of (ii).}
Fix $j \in J$ and $v \in V$. We must show that the typing constraint $E \subseteq (J \cup V) \times V$ of def \ref{def-cdmg} is compatible with the ordered pair $(j, v)$ -- i.e.\ that $(j, v) \in (J \cup V) \times V$, so that nothing in def \ref{def-cdmg} forbids the membership $(j, v) \in E$. The membership $(j, v) \in (J \cup V) \times V$ is, by the definition of a Cartesian product, the conjunction
\[ j \in J \cup V \;\land\; v \in V. \]
The second conjunct is the standing hypothesis $v \in V$. For the first conjunct, $j \in J$ gives $j \in J \cup V$ by inclusion of the left summand into the union. Hence $(j, v) \in (J \cup V) \times V$, which is the typing precondition for an element of $E$. No further restriction in def \ref{def-cdmg} excludes the pair $(j, v)$ from $E$ (the only other constraints concern $L$, not $E$, and the disjointness $J \cap V = \emptyset$ is compatible with $(j, v)$ since $j \in J$ and $v \in V$ are unrelated by intersection). Thus the inclusion $(j, v) \in E$ is admissible by the definition, which is the permission statement claimed. (We emphasise that this argument does not assert that $(j, v) \in E$ holds for any particular CDMG: it asserts only that the definition of CDMG does not preclude it.)

\medskip
\emph{Proof of (iii).}
Fix $j_1 \in J$ and $j_2 \in J$, and suppose, for contradiction, that $j_1$ and $j_2$ are adjacent in $G$. By def \ref{def-edge-relations} item~i., adjacency unfolds to the disjunction
\[ (j_1, j_2) \in E \;\lor\; (j_2, j_1) \in E \;\lor\; (j_1, j_2) \in L. \]
We rule out all three disjuncts.
\begin{itemize}
    \item If $(j_1, j_2) \in E$, then $E \subseteq (J \cup V) \times V$ (def \ref{def-cdmg}) forces the second component $j_2 \in V$. But $j_2 \in J$, so $j_2 \in J \cap V = \emptyset$ (def \ref{def-cdmg}), a contradiction.
    \item If $(j_2, j_1) \in E$, the same constraint applied to the second component yields $j_1 \in V$; combined with $j_1 \in J$ this gives $j_1 \in J \cap V = \emptyset$, a contradiction.
    \item If $(j_1, j_2) \in L$, then $L \subseteq V \times V$ (def \ref{def-cdmg}) forces both components into $V$; in particular $j_1 \in V$, hence $j_1 \in J \cap V = \emptyset$, a contradiction. (The same step also yields $j_2 \in V$; either component on its own suffices.)
\end{itemize}
All three disjuncts therefore fail, so $j_1$ and $j_2$ are not adjacent; equivalently the conjunction $(j_1, j_2) \notin E \land (j_2, j_1) \notin E \land (j_1, j_2) \notin L$ holds, as stated.

The argument above does not separate the cases $j_1 = j_2$ and $j_1 \neq j_2$: the typing-driven contradictions apply uniformly. We note in passing that the diagonal case $j_1 = j_2 =: j$ also follows: the $E$-branches give $j \in V$ as above, and the $L$-branch additionally contradicts the irreflexivity of $L$ from def \ref{def-cdmg} (which forbids $(j, j) \in L$) on top of forcing $j \in V$. No diagonal subtlety arises.
\end{proof}
\end{minipage}\hfill
\begin{minipage}[t]{0.49\textwidth}
\begin{minted}{lean4}
variable {Node : Type*} [DecidableEq Node]

theorem no_arrowhead_into_J (G : CDMG Node) {j : Node} (hj : j ∈ G.J)
    {v : Node} (hv : v ∈ G.J ∪ G.V) :
    ¬ G.hus j v

theorem J_to_V_edge_admissible (G : CDMG Node) {j : Node} (hj : j ∈ G.J)
    {v : Node} (hv : v ∈ G.V) :
    j ∈ G.J ∪ G.V ∧ v ∈ G.V

theorem J_nodes_not_adjacent (G : CDMG Node) {j₁ : Node} (hj₁ : j₁ ∈ G.J)
    {j₂ : Node} (hj₂ : j₂ ∈ G.J) :
    ¬ G.adjacent j₁ j₂
\end{minted}
\end{minipage}

\vspace{0.5em}
\noindent\textbf{Notes:}\par
% NOTES-BEGIN claim_3_1

% NOTES-END claim_3_1
\vspace{5em}
\noindent\hrulefill
\par\bigskip
%% ROW-END claim_3_1

%% ROW-BEGIN def_3_4
\subsection*{Definition: Walks}

\noindent
\begin{minipage}[t]{0.49\textwidth}
\begin{Def}[Walks]\label{def:walks}
    Let $G = (J, V, E, L)$ be a CDMG (def \ref{def-cdmg}); throughout we use the notation of def \ref{not-cdmg} and the edge classifications of def \ref{def-edge-relations}. Let $v, w \in J \cup V$.

    \begin{enumerate}[label=\roman*.)]
        \item \emph{Walk.}
            A \emph{walk} from $v$ to $w$ in $G$ is, for some integer $n \ge 0$, a tuple
            \[
                (v_0, a_0, v_1, a_1, \dots, v_{n-1}, a_{n-1}, v_n)
            \]
            consisting of $n + 1$ vertices $v_0, v_1, \dots, v_n \in J \cup V$ and $n$ ordered pairs $a_0, a_1, \dots, a_{n-1} \in (J \cup V) \times (J \cup V)$, such that:
            \begin{enumerate}[label=(\alph*)]
                \item $v_0 = v$ and $v_n = w$;
                \item for every $k \in \{0, 1, \dots, n - 1\}$,
                    \[
                        \bigl(a_k = (v_k, v_{k+1}) \;\land\; a_k \in E \cup L\bigr)
                        \;\lor\;
                        \bigl(a_k = (v_{k+1}, v_k) \;\land\; a_k \in E\bigr).
                    \]
            \end{enumerate}
            The case $n = 0$ is admitted: the \emph{trivial walk} is the single-vertex tuple $(v_0)$ with $v_0 = v = w$, no edges, and condition~(b) vacuously satisfied; it is a walk only when $v = w$. For $n \ge 1$, vertices in the tuple may repeat, i.e.\ the same node may appear multiple times among $v_0, \dots, v_n$.

            \emph{End-node classifications.}
            For a non-trivial walk ($n \ge 1$) we further classify behaviour at the two end-nodes:
            \begin{itemize}
                \item the walk is \emph{into $v_0$} iff $a_0$ is an edge into $v_0$ in the sense of def \ref{def-edge-relations}, item~ii.; concretely, iff $\bigl(a_0 = (v_1, v_0) \in E\bigr) \,\lor\, \bigl(a_0 = (v_0, v_1) \in L\bigr)$;
                \item the walk is \emph{out of $v_0$} iff $a_0$ is an edge out of $v_0$ in the sense of def \ref{def-edge-relations}, item~iii.; concretely, iff $a_0 = (v_0, v_1) \in E$;
                \item the walk is \emph{into $v_n$} iff $a_{n-1}$ is an edge into $v_n$ in the sense of def \ref{def-edge-relations}, item~ii.; concretely, iff $\bigl(a_{n-1} = (v_{n-1}, v_n) \in E\bigr) \,\lor\, \bigl(a_{n-1} = (v_{n-1}, v_n) \in L\bigr)$;
                \item the walk is \emph{out of $v_n$} iff $a_{n-1}$ is an edge out of $v_n$ in the sense of def \ref{def-edge-relations}, item~iii.; concretely, iff $a_{n-1} = (v_n, v_{n-1}) \in E$.
            \end{itemize}
            On the trivial walk ($n = 0$) all four classifications are vacuously false (neither $a_0$ nor $a_{n-1}$ exists). On a non-trivial walk with a bidirected first edge ($a_0 = (v_0, v_1) \in L$) the walk is into $v_0$ but not out of $v_0$; hence ``into $v_0$'' and ``out of $v_0$'' are mutually exclusive but not jointly exhaustive over walks (and analogously at $v_n$).

        \item \emph{Directed walk.}
            A \emph{directed walk} from $v$ to $w$ in $G$ is a walk $(v_0, a_0, v_1, \dots, a_{n-1}, v_n)$ from $v$ to $w$ (for some $n \ge 0$) that, in addition, satisfies
            \[
                a_k = (v_k, v_{k+1}) \in E \qquad \text{for every } k \in \{0, 1, \dots, n - 1\}.
            \]
            The trivial walk ($n = 0$, $v = w$) is admitted as a directed walk from $v$ to itself.

        \item \emph{Bidirected walk.}
            A \emph{bidirected walk} from $v$ to $w$ in $G$ is a walk $(v_0, a_0, v_1, \dots, a_{n-1}, v_n)$ from $v$ to $w$ (for some $n \ge 0$) that, in addition, satisfies
            \[
                a_k = (v_k, v_{k+1}) \in L \qquad \text{for every } k \in \{0, 1, \dots, n - 1\}
            \]
            (equivalently, by the symmetry of $L$ from def \ref{def-cdmg}, $(v_{k+1}, v_k) \in L$ for every such $k$). The trivial walk ($n = 0$, $v = w$) is admitted as a bidirected walk from $v$ to itself.

        \item \emph{Collider walk.}
            A \emph{collider walk} from $v$ to $w$ in $G$ is a walk $(v_0, a_0, v_1, \dots, a_{n-1}, v_n)$ from $v$ to $w$ (for some $n \ge 0$) satisfying the case-distinguished constraints below.
            \begin{itemize}
                \item \textbf{Case $n = 0$:} no further constraint; the trivial walk is a collider walk from $v$ to itself when $v = w$.
                \item \textbf{Case $n = 1$:} the single edge $a_0$ is required to be \emph{bidirected}, namely
                    \[
                        a_0 = (v_0, v_1) \in L \qquad \text{(equivalently, } (v, w) \in L\text{).}
                    \]
                    Purely directed edges $(v, w) \in E$ or $(w, v) \in E$ are \emph{not} admitted as collider walks of length~$1$.
                \item \textbf{Case $n \ge 2$:} the edges satisfy:
                    \begin{enumerate}[label=(\alph*)]
                        \item (the first edge places an arrowhead at $v_1$) $\bigl(a_0 = (v_0, v_1) \in E\bigr) \,\lor\, \bigl(a_0 = (v_0, v_1) \in L\bigr)$;
                        \item (every interior edge is bidirected) for every $k \in \{1, 2, \dots, n - 2\}$, $a_k = (v_k, v_{k+1}) \in L$;
                        \item (the last edge places an arrowhead at $v_{n-1}$) $\bigl(a_{n-1} = (v_n, v_{n-1}) \in E\bigr) \,\lor\, \bigl(a_{n-1} = (v_{n-1}, v_n) \in L\bigr)$.
                    \end{enumerate}
                    Under (a)--(c), every interior node $v_k$ (i.e.\ $1 \le k \le n - 1$) has an arrowhead pointing toward it from each of its two incident walk-edges $a_{k-1}$ and $a_k$, recovering the verbal characterisation ``every node strictly between $v$ and $w$ is a collider''.
            \end{itemize}
            \emph{Reconciliation with the source block's $n = 1$ inline note.}
            The source block contains the inline remark ``Note that for $n = 1$ this reads: $v \sus w \in G$''. This remark is \emph{overridden} by the strict reading of Case $n = 1$ above: for $n = 1$, a collider walk requires its unique edge to be bidirected, $a_0 = (v_0, v_1) \in L$. The relaxed reading ``any $v \sus w$ edge'' is \emph{not} adopted here; in particular, neither $(v, w) \in E$ nor $(w, v) \in E$ qualifies as a collider walk of length~$1$.

        \item \emph{Path.}
            A walk $(v_0, a_0, v_1, \dots, a_{n-1}, v_n)$ from $v$ to $w$ in $G$ is called a \emph{path} iff its vertex sequence contains no repetitions, i.e.\ $v_i \ne v_j$ whenever $0 \le i < j \le n$. (Equivalently, the tuple $(v_0, v_1, \dots, v_n)$ consists of $n + 1$ pairwise distinct entries. The trivial walk $n = 0$ is vacuously a path.)

        \item \emph{Bifurcation.}
            A \emph{bifurcation} between $v$ and $w$ in $G$ is a walk $(v_0, a_0, v_1, \dots, a_{n-1}, v_n)$ from $v$ to $w$ for which there exists an index $k$ with $1 \le k \le n$ (in particular $n \ge 1$) such that all of the following hold:
            \begin{enumerate}[label=(\alph*)]
                \item \emph{Distinct end-nodes appearing exactly once.}
                    $v \ne w$, and the end-nodes each occur exactly once in the vertex sequence:
                    \[
                        v_0 \notin \{v_1, v_2, \dots, v_n\}
                        \qquad\text{and}\qquad
                        v_n \notin \{v_0, v_1, \dots, v_{n-1}\}.
                    \]
                \item \emph{Left subwalk is a directed walk from $v_{k-1}$ to $v_0$.}
                    For every $j \in \{0, 1, \dots, k - 2\}$,
                    \[
                        a_j = (v_{j+1}, v_j) \in E.
                    \]
                    (When $k = 1$ the left subwalk is the trivial walk $(v_0)$ and this constraint is vacuous.)
                \item \emph{Right subwalk is a directed walk from $v_k$ to $v_n$.}
                    For every $j \in \{k, k + 1, \dots, n - 1\}$,
                    \[
                        a_j = (v_j, v_{j+1}) \in E.
                    \]
                    (When $k = n$ the right subwalk is the trivial walk $(v_n)$ and this constraint is vacuous.)
                \item \emph{Middle (hinge) edge.}
                    The edge $a_{k-1}$ between $v_{k-1}$ and $v_k$ satisfies
                    \[
                        \bigl(a_{k-1} = (v_k, v_{k-1}) \in E\bigr) \;\lor\; \bigl(a_{k-1} = (v_{k-1}, v_k) \in L\bigr).
                    \]
                \item \emph{(Addition: each end-node has exactly one arrowhead pointing toward it.)}
                    $a_0$ is an edge into $v_0$ and $a_{n-1}$ is an edge into $v_n$, in the sense of def \ref{def-edge-relations}, item~ii.; equivalently, both end-nodes $v_0$ and $v_n$ have exactly one arrowhead pointing toward them (contributed by their unique incident walk-edges $a_0$ and $a_{n-1}$ respectively, uniqueness following from clause~(a)).

                    Spelled out as a case analysis on $k$:
                    \begin{itemize}
                        \item if $k \ge 2$: clause~(b) gives $a_0 = (v_1, v_0) \in E$, which is automatically an edge into $v_0$ (no further restriction);
                        \item if $k = 1$: $a_0 = a_{k-1}$ \emph{is} the middle edge; both alternatives admitted by~(d) — namely $a_0 = (v_1, v_0) \in E$ and $a_0 = (v_0, v_1) \in L$ — are edges into $v_0$, so no further restriction at $v_0$;
                        \item if $k \le n - 1$: clause~(c) gives $a_{n-1} = (v_{n-1}, v_n) \in E$, which is automatically an edge into $v_n$ (no further restriction);
                        \item if $k = n$: $a_{n-1} = a_{k-1}$ \emph{is} the middle edge; among the two alternatives admitted by~(d), only the bidirected alternative $a_{n-1} = (v_{n-1}, v_n) \in L$ is an edge into $v_n$, whereas the directed alternative $a_{n-1} = (v_n, v_{n-1}) \in E$ places its arrowhead at $v_{n-1}$ rather than $v_n$ and is therefore \emph{excluded}.
                    \end{itemize}
                    In particular, the degenerate configuration ``$k = n$ together with the directed hinge $a_{n-1} = (v_n, v_{n-1}) \in E$'' — which would otherwise admit a tuple of the form $v_n \to v_{n-1} \to \cdots \to v_1 \to v_0$ (a plain directed walk from $v_n$ to $v_0$) as a ``bifurcation'' — is \emph{not} a bifurcation under the present definition.
            \end{enumerate}
            \emph{Source of the bifurcation.}
            If clause~(d) is realised by the \emph{directed} alternative $a_{k-1} = (v_k, v_{k-1}) \in E$, then we call $v_k$ the \emph{source} of the bifurcation. If clause~(d) is realised by the \emph{bidirected} alternative $a_{k-1} = (v_{k-1}, v_k) \in L$, no source is defined. Combined with clause~(e), whenever a source $v_k$ is defined we automatically have $1 \le k \le n - 1$ (the case $k = n$ with directed hinge is excluded by~(e)), so the source $v_k$ is an interior vertex of the walk and is distinct from both end-nodes $v_0 = v$ and $v_n = w$.
    \end{enumerate}
\end{Def}
\end{minipage}\hfill
\begin{minipage}[t]{0.49\textwidth}
\begin{minted}{lean4}
variable {Node : Type*} [DecidableEq Node]

def WalkStep (G : CDMG Node) (u : Node) (a : Node × Node) (v : Node) : Prop :=
  (a = (u, v) ∧ (a ∈ G.E ∨ a ∈ G.L)) ∨ (a = (v, u) ∧ a ∈ G.E)

variable {G : CDMG Node}

def vertices : ∀ {u v : Node}, Walk G u v → List Node
  | _, _, .nil v _ => [v]
  | u, _, .cons _ _ _ p => u :: p.vertices

def edges : ∀ {u v : Node}, Walk G u v → List (Node × Node)
  | _, _, .nil _ _ => []
  | _, _, .cons _ a _ p => a :: p.edges

def IsColliderRest : ∀ {u v : Node}, Walk G u v → Prop
  | _, _, .nil _ _ => True
  | u, _, .cons v a _ (.nil _ _) =>
      (a = (v, u) ∧ a ∈ G.E) ∨ (a = (u, v) ∧ a ∈ G.L)
  | u, _, .cons v a _ (p@(.cons _ _ _ _)) =>
      a = (u, v) ∧ a ∈ G.L ∧ p.IsColliderRest

def IsBifurcationWithSplit : ∀ {u v : Node}, Walk G u v → ℕ → Prop
  | _, _, .nil _ _, _ => False
  | u, _, .cons v a _ (.nil _ _), 0 => a = (u, v) ∧ a ∈ G.L
  | u, _, .cons v a _ (p@(.cons _ _ _ _)), 0 =>
      ((a = (v, u) ∧ a ∈ G.E) ∨ (a = (u, v) ∧ a ∈ G.L)) ∧ p.IsDirectedWalk
  | u, _, .cons v a _ p, k + 1 =>
      a = (v, u) ∧ a ∈ G.E ∧ p.IsBifurcationWithSplit k

def IsBifurcationDirectedHingeWithSplit : ∀ {u v : Node}, Walk G u v → ℕ → Prop
  | _, _, .nil _ _, _ => False
  | _, _, .cons _ _ _ (.nil _ _), 0 => False
  | u, _, .cons v a _ (p@(.cons _ _ _ _)), 0 =>
      a = (v, u) ∧ a ∈ G.E ∧ p.IsDirectedWalk
  | u, _, .cons v a _ p, k + 1 =>
      a = (v, u) ∧ a ∈ G.E ∧ p.IsBifurcationDirectedHingeWithSplit k

inductive Walk (G : CDMG Node) : Node → Node → Type _ where
  | nil (v : Node) (hv : v ∈ G) : Walk G v v
  | cons {u w : Node} (v : Node) (a : Node × Node)
      (h : G.WalkStep u a v) (p : Walk G v w) : Walk G u w

def intoStart : ∀ {u v : Node}, Walk G u v → Prop
  | _, _, .nil _ _ => False
  | u, _, .cons _ a _ _ => G.into u a

def outOfStart : ∀ {u v : Node}, Walk G u v → Prop
  | _, _, .nil _ _ => False
  | u, _, .cons _ a _ _ => G.outOf u a

def intoEnd {u v : Node} (p : Walk G u v) : Prop :=
  match p.edges.getLast? with
  | none => False
  | some a => G.into v a

def outOfEnd {u v : Node} (p : Walk G u v) : Prop :=
  match p.edges.getLast? with
  | none => False
  | some a => G.outOf v a

def IsDirectedWalk : ∀ {u v : Node}, Walk G u v → Prop
  | _, _, .nil _ _ => True
  | u, _, .cons v a _ p => a = (u, v) ∧ a ∈ G.E ∧ p.IsDirectedWalk

def IsBidirectedWalk : ∀ {u v : Node}, Walk G u v → Prop
  | _, _, .nil _ _ => True
  | u, _, .cons v a _ p => a = (u, v) ∧ a ∈ G.L ∧ p.IsBidirectedWalk

def IsColliderWalk : ∀ {u v : Node}, Walk G u v → Prop
  | _, _, .nil _ _ => True
  | u, _, .cons v a _ (.nil _ _) => a = (u, v) ∧ a ∈ G.L
  | u, _, .cons v a _ (p@(.cons _ _ _ _)) =>
      (a = (u, v) ∧ (a ∈ G.E ∨ a ∈ G.L)) ∧ p.IsColliderRest

def IsPath {u v : Node} (p : Walk G u v) : Prop := p.vertices.Nodup

def IsBifurcation {u v : Node} (p : Walk G u v) : Prop :=
  u ≠ v ∧
  u ∉ p.vertices.tail ∧
  v ∉ p.vertices.dropLast ∧
  ∃ i, p.IsBifurcationWithSplit i

def IsBifurcationSource {u v : Node} (p : Walk G u v) (x : Node) : Prop :=
  u ≠ v ∧
  u ∉ p.vertices.tail ∧
  v ∉ p.vertices.dropLast ∧
  ∃ i, p.IsBifurcationDirectedHingeWithSplit i ∧ p.vertices[i + 1]? = some x
\end{minted}
\end{minipage}

\vspace{0.5em}
\noindent\textbf{Notes:}\par
% NOTES-BEGIN def_3_4

% NOTES-END def_3_4
\vspace{5em}
\noindent\hrulefill
\par\bigskip
%% ROW-END def_3_4

%% ROW-BEGIN def_3_5
\subsection*{Definition: Family relationships}

\noindent
\begin{minipage}[t]{0.49\textwidth}
\begin{Def}[Family relationships]
    Let $G = (J, V, E, L)$ be a CDMG (def \ref{def-cdmg}); throughout we use the notation of def \ref{not-cdmg} and the walk definitions of def \ref{def:walks}. Let $v, w \in J \cup V$ and let $A \ins J \cup V$ be a subset of nodes. We adopt the convention that each individual form $\Pa^G(v)$, $\Ch^G(v)$, $\Sib^G(v)$, $\Anc^G(v)$, $\Desc^G(v)$, $\Sc^G(v)$, $\Dist^G(v)$ defined below is defined for any $v \in J \cup V$ (not only $v \in V$); in particular, input nodes $v \in J$ admit the same set-builder definitions, and each set form $\Pa^G(A)$, $\Ch^G(A)$, $\Anc^G(A)$, $\Desc^G(A)$, $\Sc^G(A)$, $\Dist^G(A)$ is well-typed for $A \ins J \cup V$ because every summand $\Pa^G(v)$, $\Ch^G(v)$, $\Anc^G(v)$, $\Desc^G(v)$, $\Sc^G(v)$, $\Dist^G(v)$ is defined.

    \begin{enumerate}[label=\roman*.)]
        \item \emph{Parents.}
            The set of \emph{parents} of $v$ in $G$ is
            \[
                \Pa^G(v) := \lC w \in J \cup V \,|\, (w, v) \in E \rC.
            \]
            The set of \emph{parents} of $A$ in $G$ is
            \[
                \Pa^G(A) := \bigcup_{v \in A} \Pa^G(v).
            \]

        \item \emph{Children.}
            The set of \emph{children} of $v$ in $G$ is
            \[
                \Ch^G(v) := \lC w \in J \cup V \,|\, (v, w) \in E \rC.
            \]
            The set of \emph{children} of $A$ in $G$ is
            \[
                \Ch^G(A) := \bigcup_{v \in A} \Ch^G(v).
            \]

        \item \emph{Siblings.}
            The set of \emph{siblings} of $v$ in $G$ is
            \[
                \Sib^G(v) := \lC w \in J \cup V \,|\, (v, w) \in L \rC.
            \]

        \item \emph{Ancestors.}
            The set of \emph{ancestors} of $v$ in $G$ is
            \[
                \Anc^G(v) := \lC w \in J \cup V \,|\, \text{there exists a directed walk from } w \text{ to } v \text{ in } G \rC,
            \]
            where ``directed walk'' is in the sense of def \ref{def:walks}, item~ii., admitted to be of any length $n \ge 0$. In particular, the length-$0$ (trivial) directed walk at $v$ (admitted by def \ref{def:walks}, item~ii., as the single-vertex tuple $(v)$ with no edges) witnesses $v \in \Anc^G(v)$. Hence
            \[
                v \in \Anc^G(v) \qquad \text{for every } v \in J \cup V,
            \]
            \emph{unconditionally}: this holds regardless of whether any edge (in particular any directed self-loop) incident to $v$ exists in $G$.

            The set of \emph{ancestors} of $A$ in $G$ is
            \[
                \Anc^G(A) := \bigcup_{v \in A} \Anc^G(v),
            \]
            and consequently $A \ins \Anc^G(A)$ unconditionally.

        \item \emph{Descendants.}
            The set of \emph{descendants} of $v$ in $G$ is
            \[
                \Desc^G(v) := \lC w \in J \cup V \,|\, \text{there exists a directed walk from } v \text{ to } w \text{ in } G \rC,
            \]
            with ``directed walk'' as in def \ref{def:walks}, item~ii., admitted to be of any length $n \ge 0$. The length-$0$ (trivial) directed walk at $v$ witnesses $v \in \Desc^G(v)$, so
            \[
                v \in \Desc^G(v) \qquad \text{for every } v \in J \cup V,
            \]
            \emph{unconditionally} (regardless of whether any edge, in particular any directed self-loop, incident to $v$ exists in $G$).

            The set of \emph{descendants} of $A$ in $G$ is
            \[
                \Desc^G(A) := \bigcup_{v \in A} \Desc^G(v),
            \]
            and consequently $A \ins \Desc^G(A)$ unconditionally.

        \item \emph{Non-descendants.}
            The set of \emph{non-descendants} of $A$ in $G$ is
            \[
                \NonDesc^G(A) := (J \cup V) \sm \Desc^G(A).
            \]

        \item \emph{Strongly connected component.}
            The \emph{strongly connected component} of $v$ in $G$ is
            \[
                \Sc^G(v) := \Anc^G(v) \cap \Desc^G(v).
            \]
            Since $v \in \Anc^G(v)$ and $v \in \Desc^G(v)$ unconditionally (items~iv.\ and~v.), it follows that
            \[
                v \in \Sc^G(v) \qquad \text{for every } v \in J \cup V,
            \]
            \emph{unconditionally} (regardless of whether any edge incident to $v$ exists in $G$).

            The (union of) \emph{strongly connected components} of $A$ in $G$ is
            \[
                \Sc^G(A) := \bigcup_{v \in A} \Sc^G(v),
            \]
            and consequently $A \ins \Sc^G(A)$ unconditionally.

        \item \emph{District.}
            The \emph{district} of $v$ in $G$ is
            \[
                \Dist^G(v) := \lC w \in J \cup V \,|\, \text{there exists a bidirected walk from } v \text{ to } w \text{ in } G \rC,
            \]
            where ``bidirected walk'' is in the sense of def \ref{def:walks}, item~iii., admitted to be of any length $n \ge 0$. The lecture notes' displayed shape ``$v \huh v_1 \huh \cdots \huh v_{n-1} \huh w$'' is to be read here purely as syntactic sugar for ``walk all of whose edges lie in $L$''; the intermediate-node indexing $v_1, \dots, v_{n-1}$ imposes no lower bound on $n$, and in particular admits $n = 0$. The length-$0$ (trivial) bidirected walk at $v$ witnesses $v \in \Dist^G(v)$, so
            \[
                v \in \Dist^G(v) \qquad \text{for every } v \in J \cup V,
            \]
            \emph{unconditionally}: neither a bidirected self-loop at $v$ nor a length-$2$ bidirected cycle through $v$ is required for self-membership (and indeed a bidirected self-loop at $v$ is excluded a priori by the irreflexivity of $L$ in def \ref{def-cdmg}).

            The \emph{district} of $A$ in $G$ is
            \[
                \Dist^G(A) := \bigcup_{v \in A} \Dist^G(v),
            \]
            and consequently $A \ins \Dist^G(A)$ unconditionally.
    \end{enumerate}
\end{Def}
\end{minipage}\hfill
\begin{minipage}[t]{0.49\textwidth}
\begin{minted}{lean4}
variable {Node : Type*} [DecidableEq Node]

def Pa (G : CDMG Node) (v : Node) : Set Node :=
  {w | w ∈ G ∧ (w, v) ∈ G.E}

def PaSet (G : CDMG Node) (A : Set Node) : Set Node :=
  ⋃ v ∈ A, G.Pa v

def Ch (G : CDMG Node) (v : Node) : Set Node :=
  {w | w ∈ G ∧ (v, w) ∈ G.E}

def ChSet (G : CDMG Node) (A : Set Node) : Set Node :=
  ⋃ v ∈ A, G.Ch v

def Sib (G : CDMG Node) (v : Node) : Set Node :=
  {w | w ∈ G ∧ (v, w) ∈ G.L}

def Anc (G : CDMG Node) (v : Node) : Set Node :=
  {w | w ∈ G ∧ ∃ p : Walk G w v, p.IsDirectedWalk}

def AncSet (G : CDMG Node) (A : Set Node) : Set Node :=
  ⋃ v ∈ A, G.Anc v

def Desc (G : CDMG Node) (v : Node) : Set Node :=
  {w | w ∈ G ∧ ∃ p : Walk G v w, p.IsDirectedWalk}

def DescSet (G : CDMG Node) (A : Set Node) : Set Node :=
  ⋃ v ∈ A, G.Desc v

def NonDesc (G : CDMG Node) (A : Set Node) : Set Node :=
  ((G.J ∪ G.V : Finset Node) : Set Node) \ G.DescSet A

def Sc (G : CDMG Node) (v : Node) : Set Node :=
  G.Anc v ∩ G.Desc v

def ScSet (G : CDMG Node) (A : Set Node) : Set Node :=
  ⋃ v ∈ A, G.Sc v

def Dist (G : CDMG Node) (v : Node) : Set Node :=
  {w | w ∈ G ∧ ∃ p : Walk G v w, p.IsBidirectedWalk}

def DistSet (G : CDMG Node) (A : Set Node) : Set Node :=
  ⋃ v ∈ A, G.Dist v
\end{minted}
\end{minipage}

\vspace{0.5em}
\noindent\textbf{Notes:}\par
% NOTES-BEGIN def_3_5

% NOTES-END def_3_5
\vspace{5em}
\noindent\hrulefill
\par\bigskip
%% ROW-END def_3_5

%% ROW-BEGIN def_3_6
\subsection*{Definition: Acyclicity}

\noindent
\begin{minipage}[t]{0.49\textwidth}
\begin{Def}[Acyclicity]
    \label{def-acylic}
    Let $G = (J, V, E, L)$ be a CDMG (def \ref{def-cdmg}); throughout we use the notation of def \ref{not-cdmg} and the walk definitions of def \ref{def:walks}. By a \emph{directed walk from $v$ to $w$ in $G$} (for $v, w \in J \cup V$) we mean a directed walk in the sense of def \ref{def:walks}, item~ii., admitted to be of any length $n \ge 0$: a tuple $(v_0, a_0, v_1, a_1, \dots, a_{n-1}, v_n)$ with $v_0, v_1, \dots, v_n \in J \cup V$, $v_0 = v$, $v_n = w$, and $a_k = (v_k, v_{k+1}) \in E$ for every $k \in \{0, 1, \dots, n - 1\}$ (in the LN's diagrammatic shorthand of def \ref{not-cdmg}, item~2: $v = v_0 \tuh v_1 \tuh \cdots \tuh v_{n-1} \tuh v_n = w$). Such a directed walk is called \emph{non-trivial} iff its length $n$ satisfies $n \ge 1$, equivalently iff it traverses at least one edge; the length-$0$ \emph{trivial} directed walk $(v_0)$ at $v_0 = v = w$ (admitted by def \ref{def:walks}, item~ii., as the single-vertex tuple with no edges) is excluded.

    The CDMG $G$ is called \emph{acyclic} iff
    \[
        \text{for every } v \in J \cup V, \quad \text{there does not exist any non-trivial directed walk from } v \text{ to itself in } G.
    \]
    Spelled out as a first-order statement over the components $J, V, E$ of $G$: for every $v \in J \cup V$, every integer $n \ge 1$, and every tuple $(v_0, v_1, \dots, v_n) \in (J \cup V)^{n+1}$ with $v_0 = v_n = v$, there exists some $k \in \{0, 1, \dots, n - 1\}$ with $(v_k, v_{k+1}) \notin E$ (i.e.\ no such vertex tuple is the vertex sequence of a directed walk from $v$ to itself in $G$).

    \emph{Consequence: no directed self-loops on output nodes.}
    For any $v \in V$, a directed self-loop on $v$ -- i.e.\ the ordered pair $(v, v) \in E$, equivalently $v \tuh v \in G$ in the notation of def \ref{not-cdmg}, item~2 -- yields the length-$1$ directed walk $\bigl(v, (v, v), v\bigr)$ from $v$ to itself, which is non-trivial in the sense above. Hence if $G$ is acyclic, then $(v, v) \notin E$ for every $v \in V$. (For $v \in J$, the CDMG-type constraint $E \subseteq (J \cup V) \times V$ together with $J \cap V = \emptyset$ already excludes $(v, v) \in E$ a priori, so the acyclicity hypothesis adds no further constraint at $v \in J$.) In particular, an acyclic CDMG contains no directed self-loops on any $v \in V$.
\end{Def}
\end{minipage}\hfill
\begin{minipage}[t]{0.49\textwidth}
\begin{minted}{lean4}
variable {Node : Type*} [DecidableEq Node]

def IsAcyclic (G : CDMG Node) : Prop :=
  ∀ v ∈ G, ¬ ∃ p : Walk G v v, p.IsDirectedWalk ∧ p.length ≥ 1
\end{minted}
\end{minipage}

\vspace{0.5em}
\noindent\textbf{Notes:}\par
% NOTES-BEGIN def_3_6

% NOTES-END def_3_6
\vspace{5em}
\noindent\hrulefill
\par\bigskip
%% ROW-END def_3_6

%% ROW-BEGIN def_3_7
\subsection*{Definition: CDMGTypes}

\noindent
\begin{minipage}[t]{0.49\textwidth}
\begin{Def}[CDMGTypes]
    \label{def-cdmg-types}
    Let $G = (J, V, E, L)$ be a CDMG in the sense of def \ref{def-cdmg}. Items~i.--vii.\ below introduce seven names; each name is an \emph{attribute} that $G$ may bear, declared by a clause of the form ``$G$ is called \emph{<name>} iff <condition>''. Every condition is a conjunction drawn from the three atomic conditions
    \begin{itemize}
        \item ``$G$ is acyclic'', the acyclicity predicate of def \ref{def-acylic};
        \item ``$J = \emptyset$'', read as equality of the (finite) input-node set $J$ with the empty set, as elements of the finite-subset type of the underlying node universe of $G$;
        \item ``$L = \emptyset$'', read as equality of the (finite) bidirected-edge set $L$ with the empty set, as elements of the finite-subset type of $V \times V$ (the ambient type on which $L$ lives per def \ref{def-cdmg}).
    \end{itemize}

    \emph{Nested-name semantics.} The seven names below do \emph{not} partition CDMGs. ``$G$ is called \emph{<name>}'' is read as ``$G$ \emph{satisfies the attribute} \emph{<name>}'', and a single CDMG $G$ may simultaneously satisfy several of the listed conditions, in which case it simultaneously bears each corresponding name. In particular, whenever the condition of item~vii.\ (DAG) holds, so do the conditions of items~i.--vi.\ and the base attribute ``CDMG'' of def \ref{def-cdmg}, so $G$ then bears all eight names at once.

    \begin{enumerate}[label=\roman*.)]
        \item \emph{Conditional Acyclic Directed Mixed Graph (CADMG).}
              $G$ is called a \emph{CADMG} iff $G$ is acyclic in the sense of def \ref{def-acylic}.

        \item \emph{Directed Mixed Graph (DMG).}
              $G$ is called a \emph{DMG} iff $J = \emptyset$.

        \item \emph{Acyclic Directed Mixed Graph (ADMG).}
              $G$ is called an \emph{ADMG} iff $G$ is acyclic in the sense of def \ref{def-acylic} \emph{and} $J = \emptyset$.

        \item \emph{Conditional Directed Graph (CDG).}
              $G$ is called a \emph{CDG} iff $L = \emptyset$.

        \item \emph{Directed Graph (DG).}
              $G$ is called a \emph{DG} iff $J = \emptyset$ \emph{and} $L = \emptyset$.

        \item \emph{Conditional Directed Acyclic Graph (CDAG).}
              $G$ is called a \emph{CDAG} iff $G$ is acyclic in the sense of def \ref{def-acylic} \emph{and} $L = \emptyset$.

        \item \emph{Directed Acyclic Graph (DAG).}
              $G$ is called a \emph{DAG} iff $G$ is acyclic in the sense of def \ref{def-acylic}, $J = \emptyset$, \emph{and} $L = \emptyset$.
    \end{enumerate}

    \emph{Coverage.} The three atomic conditions ``$G$ is acyclic'', ``$J = \emptyset$'', ``$L = \emptyset$'' are independent boolean predicates on $G$, yielding $2^3 = 8$ possible truth-value tuples. Each of items~i.--vii.\ requires a distinct non-empty subset of these three to hold (CADMG: $\{\text{acyclic}\}$; DMG: $\{J = \emptyset\}$; ADMG: $\{\text{acyclic}, J = \emptyset\}$; CDG: $\{L = \emptyset\}$; DG: $\{J = \emptyset, L = \emptyset\}$; CDAG: $\{\text{acyclic}, L = \emptyset\}$; DAG: $\{\text{acyclic}, J = \emptyset, L = \emptyset\}$); the bare CDMG of def \ref{def-cdmg} (requiring the empty subset, i.e.\ none of the three) accounts for the remaining combination, so the seven named attributes together with the base CDMG attribute exhaust all $2^3 = 8$ boolean combinations. No item constrains the output-node set $V$ or the directed-edge set $E$ beyond what def \ref{def-cdmg} already requires.
\end{Def}
\end{minipage}\hfill
\begin{minipage}[t]{0.49\textwidth}
\begin{minted}{lean4}
variable {Node : Type*} [DecidableEq Node]

def IsCADMG (G : CDMG Node) : Prop := G.IsAcyclic

def IsDMG (G : CDMG Node) : Prop := G.J = ∅

def IsADMG (G : CDMG Node) : Prop := G.IsAcyclic ∧ G.J = ∅

def IsCDG (G : CDMG Node) : Prop := G.L = ∅

def IsDG (G : CDMG Node) : Prop := G.J = ∅ ∧ G.L = ∅

def IsCDAG (G : CDMG Node) : Prop := G.IsAcyclic ∧ G.L = ∅

def IsDAG (G : CDMG Node) : Prop := G.IsAcyclic ∧ G.J = ∅ ∧ G.L = ∅
\end{minted}
\end{minipage}

\vspace{0.5em}
\noindent\textbf{Notes:}\par
% NOTES-BEGIN def_3_7

% NOTES-END def_3_7
\vspace{5em}
\noindent\hrulefill
\par\bigskip
%% ROW-END def_3_7

%% ROW-BEGIN def_3_8
\subsection*{Definition: Topological order}

\noindent
\begin{minipage}[t]{0.49\textwidth}
\begin{Def}[Topological order]
    \label{def-topological-order}
    Let $G = (J, V, E, L)$ be a CDMG (def \ref{def-cdmg}); throughout we use the notation of def \ref{not-cdmg} and the family-relationship notation of def \ref{def_3_5}. By def \ref{def-cdmg}, items~i.--ii., the node set $J \cup V$ is finite (as a union of two finite sets).

    A \emph{topological order} of $G$ is a \emph{strict} total order $<$ on $J \cup V$ -- i.e.\ a binary relation $<$ on $J \cup V$ that is irreflexive, transitive, and satisfies trichotomy ($v < w$, $v = w$, or $w < v$ for every $v, w \in J \cup V$); the lecture notes' order symbol $<$ is read strictly throughout -- such that
    \[
        \forall\, v, w \in J \cup V :\quad v \in \Pa^G(w) \;\Longrightarrow\; v < w.
    \]
    The displayed implication is well-typed because, by the set-builder defining clause $\Pa^G(w) := \lC u \in J \cup V \,|\, (u, w) \in E \rC$ of def \ref{def_3_5}, item~i., any $v \in \Pa^G(w)$ automatically lies in $J \cup V$, so the comparison $v < w$ is defined.

    \emph{Equivalent indexing form.}
    Set $K := |J \cup V|$, the cardinality of the (finite) node set. Because $<$ is a strict total order on the finite set $J \cup V$, its order type is the finite ordinal $K$, so $<$ may equivalently be presented as a bijection $\{1, 2, \dots, K\} \to J \cup V$, $i \mapsto v_i$, with $v_1 < v_2 < \cdots < v_K$ under $<$. Under this presentation, the parents-precede-children condition above becomes
    \[
        \forall\, i, j \in \{1, 2, \dots, K\} :\quad v_i \in \Pa^G(v_j) \;\Longrightarrow\; i < j,
    \]
    i.e.\ \emph{parents always precede their children} in the enumeration $v_1, v_2, \dots, v_K$. The equivalence between the primary (strict-total-order) form and this indexed form relies essentially on the finiteness of $J \cup V$ guaranteed by def \ref{def-cdmg}: without finiteness, the indexed form would be strictly stronger than the primary form (it forces order type at most $\omega$).
\end{Def}
\end{minipage}\hfill
\begin{minipage}[t]{0.49\textwidth}
\begin{minted}{lean4}
variable {Node : Type*} [DecidableEq Node]

def IsTotalOrder (G : CDMG Node) (lt : Node → Node → Prop) : Prop :=
  (∀ v ∈ G, ¬ lt v v) ∧
  (∀ u ∈ G, ∀ v ∈ G, ∀ w ∈ G, lt u v → lt v w → lt u w) ∧
  (∀ v ∈ G, ∀ w ∈ G, lt v w ∨ v = w ∨ lt w v)

def IsTopologicalOrder (G : CDMG Node) (lt : Node → Node → Prop) : Prop :=
  G.IsTotalOrder lt ∧ (∀ v w, v ∈ G.Pa w → lt v w)
\end{minted}
\end{minipage}

\vspace{0.5em}
\noindent\textbf{Notes:}\par
% NOTES-BEGIN def_3_8

% NOTES-END def_3_8
\vspace{5em}
\noindent\hrulefill
\par\bigskip
%% ROW-END def_3_8

%% ROW-BEGIN claim_3_2
\subsection*{Claim: AcyclicIffTopologicalOrder}

\noindent
\begin{minipage}[t]{0.49\textwidth}
% --- statement (restated) ---------------------------------------------------
\begin{Lem}[AcyclicIffTopologicalOrder]
    Let $G = (J, V, E, L)$ be a CDMG (def \ref{def-cdmg}). Then $G$ is acyclic (in the sense of def \ref{def-acylic}) if and only if there exists a topological order of $G$ (in the sense of def \ref{def-topological-order}); equivalently, $G$ is acyclic if and only if there exists a strict total order $<$ on $J \cup V$ -- i.e.\ a binary relation $<$ on $J \cup V$ that is irreflexive, transitive, and trichotomous -- such that, for every $v, w \in J \cup V$, $v \in \Pa^G(w) \implies v < w$ (using the parent-set notation $\Pa^G$ of def \ref{def_3_5}, item~i.). The label set $L$ plays no role on either side of the biconditional.
\end{Lem}

% --- proof ------------------------------------------------------------------
% Proof lifted verbatim from the LN at lecture-notes/lecture_notes/graphs.tex:228-244.
\begin{proof}
($\Leftarrow$)
Suppose $<$ is a topological order of $G$, i.e.\ $v \in \Pa^G(w) \implies v < w$.
If there were a non-trivial directed walk $v = v_0 \tuh v_1 \tuh \cdots \tuh v_n = v$ ($n \ge 1$),
then $v_0 < v_1 < \cdots < v_n = v_0$, a contradiction.

\medskip\noindent($\Rightarrow$)
Suppose $G$ is acyclic. We construct a topological order by repeatedly selecting a node with no parents in the remaining graph.
Define a sequence $v_1, \dots, v_K$ ($K = |J \cup V|$) inductively: given $v_1,\dots,v_{i-1}$, let $R_i := (J \cup V) \sm \{v_1,\dots,v_{i-1}\}$
and $G_i$ the subgraph of $G$ induced on $R_i$.
Since $G_i$ is acyclic (as an induced subgraph of an acyclic graph) and finite, it has a node $v_i$ with $\Pa^{G_i}(v_i) = \emptyset$
(otherwise, following parent edges would produce a directed cycle).
Choose any such $v_i$.

This gives a total order $v_1 < v_2 < \cdots < v_K$.
If $v = v_j \in \Pa^G(w = v_i)$, then when $w$ was selected, $v$ was still in $R_i$, so $v$ would have been a parent of $w$ in $G_i$, contradicting $\Pa^{G_i}(v_i) = \emptyset$---unless $v$ was already selected, i.e.\ $j < i$. Hence $v < w$.
\end{proof}
\end{minipage}\hfill
\begin{minipage}[t]{0.49\textwidth}
\begin{minted}{lean4}
variable {Node : Type*} [DecidableEq Node]

theorem acyclic_iff_topological_order (G : CDMG Node) :
    G.IsAcyclic ↔ ∃ lt : Node → Node → Prop, G.IsTopologicalOrder lt
\end{minted}
\end{minipage}

\vspace{0.5em}
\noindent\textbf{Notes:}\par
% NOTES-BEGIN claim_3_2

% NOTES-END claim_3_2
\vspace{5em}
\noindent\hrulefill
\par\bigskip
%% ROW-END claim_3_2

%% ROW-BEGIN def_3_9
\subsection*{Definition: Predecessors}

\noindent
\begin{minipage}[t]{0.49\textwidth}
\begin{Def}[Predecessors]
    \label{def-predecessors}
    Let $G = (J, V, E, L)$ be a CDMG (def \ref{def-cdmg}); throughout we use the notation of def \ref{not-cdmg}. Let $<$ be a total order of $J \cup V$, read as a \emph{strict} total order in the sense of def \ref{def-topological-order} (irreflexive, transitive, and trichotomous on $J \cup V$); the symbol $<$ is read strictly throughout. The element $v$ is left unconstrained by the LN -- in particular, the LN does \emph{not} impose $v \in J \cup V$ -- and we follow that literal stance below; note that the comparison $w < v$ is only directly given by $<$ when $v \in J \cup V$, since $<$ is supplied only as a total order of $J \cup V$. Adopting the convention of def \ref{not-cdmg}, item~1, that ``$w \in G$'' abbreviates ``$w \in J \cup V$'', we write the set-builder bodies below explicitly with $w \in J \cup V$.

    The set of \emph{predecessors} of $v$ in $G$ is
    \[
        \Pred^G_<(v) := \lC w \in J \cup V \,|\, w < v \rC.
    \]
    We also put
    \[
        \Pred^G_\le(v) := \lC w \in J \cup V \,|\, w < v \rC \cup \{v\}.
    \]

    \emph{Notational quirk: the ``$\le$'' subscript vs.\ the strict body.} The subscript ``$\le$'' on $\Pred^G_\le$ marks the non-strict variant, but the set-builder body the LN writes is \emph{strict} (the membership condition is $w < v$, not $w \le v$), with the singleton $\{v\}$ adjoined separately. We implement this literal LN body $\Pred^G_<(v) \cup \{v\}$, \emph{not} the alternative $\lC w \in J \cup V \,|\, w \le v \rC$; the two forms coincide whenever $v \in J \cup V$ (irreflexivity of $<$ keeps $v$ out of the strict body, while the reflexive-closure form $w \le v$ would pick it up via $v \le v$), and they diverge only in the corner case noted next.

    \emph{Corner case $v \notin J \cup V$.} Under the literal LN body, $w < v$ for $w \in J \cup V$ is not provided by the total order $<$ (whose domain is $J \cup V$), so the natural reading is $\Pred^G_<(v) = \emptyset$; the non-strict body's adjoined singleton then degenerately gives $\Pred^G_\le(v) = \{v\}$, so $v$ is its own non-strict predecessor even when $v$ lies outside the underlying node set of $G$. We follow the LN's lack of hypothesis on $v$ and do \emph{not} impose $v \in J \cup V$ here; downstream consumers that pattern-match on the shape of an element of $\Pred^G_<(v)$ or $\Pred^G_\le(v)$ may assume $v \in J \cup V$ at the point of use as a separate hypothesis.
\end{Def}
\end{minipage}\hfill
\begin{minipage}[t]{0.49\textwidth}
\begin{minted}{lean4}
variable {Node : Type*} [DecidableEq Node]

def Pred (G : CDMG Node) (lt : Node → Node → Prop)
    (_h : G.IsTotalOrder lt) (v : Node) : Set Node :=
  {w | w ∈ G ∧ lt w v}

def PredLE (G : CDMG Node) (lt : Node → Node → Prop)
    (h : G.IsTotalOrder lt) (v : Node) : Set Node :=
  G.Pred lt h v ∪ {v}
\end{minted}
\end{minipage}

\vspace{0.5em}
\noindent\textbf{Notes:}\par
% NOTES-BEGIN def_3_9

% NOTES-END def_3_9
\vspace{5em}
\noindent\hrulefill
\par\bigskip
%% ROW-END def_3_9
